%% notes-tex/database/database-expansion-200/sections/synergies.tex
%% [[experiments.database.expansion-200]]
%% https://github.com/Mjvolk3/torchcell/tree/main/notes-tex/database/database-expansion-200/sections/synergies.tex
%% SOURCE: experiments/database/scripts/build_candidate_datasets_table.py for the table;
%% surrounding prose is hand-authored

\section{What each candidate joins to\secstatus{todo}}\label{sec:synergies}

A dataset's value in a substrate is not only what it measures. It is also what becomes
measurable once it sits beside what is already held. Table~\ref{tab:synergies} names 89
such pairs across the sixty rows of Table~\ref{tab:final}, 53 of the sixty carrying at
least one, each with the axis the two share.

\subsection{A join needs a shared axis, not a shared theme\secstatus{todo}}

The \term{join key} column is required rather than decorative. Two metabolomics papers are
not joinable because both measure metabolites; they are joinable when they measure
overlapping analytes on overlapping strains. The keys in use are the deletion collection,
the sequenced 1,011-isolate panel, a segregant panel from one cross, a guide library over
one gene set, a media formulation, a compound set, and a single gene. A pair with none of
these is a theme and is not listed.

Fifty-five of the 89 pairs, on 40 of the sixty rows, have a partner that is already
served, so that join becomes runnable on the day the candidate is ingested rather than
after two ingestions. Thirteen of those partners became served since the previous pass,
which is why pairs that were candidate-to-candidate then, Teyssonniere 2024 against
Muenzner 2024 or Hughes 2000 against Hu 2007, are now one ingestion from running.

Two kinds of key are distinguished in the table and the difference decides what a join can
do. A \emph{strain-level} key means both layers were measured on the same genotypes, so
the join is a paired observation and the residual between the layers is per strain. A
\emph{gene-level} key means one side is a single genome-wide coefficient, a half-life or a
translation rate, that applies to every strain alike. A gene-level layer cannot explain why
two strains differ; it explains why two genes differ, and the two are not interchangeable.

\subsection{The joins that motivate the ordering\secstatus{todo}}

Four groups carry most of the weight, and each is a specific question rather than an
accumulation.

\notesec{State on one genome axis} The eight rows of the first band are joins by
construction. Albert 2018 and Jakobson 2025 pair transcript and protein on one segregant
class, and Albert 2018 pairs with the served Boocock 2025 on the 393 pooled segregants
the two share. Teyssonniere 2024 pairs with the served Caudal 2024 strain by strain, 942
against 943 isolates, and with the served Muenzner 2024 on the protein layer, where the
pairing is also the de-duplication question. Grossbach 2022 and Foss 2007 are the same
cross fifteen years apart. Smith 2008 is the segregant transcriptome with a condition
axis, so it pairs with every other BY x RM row on genotype and with Jackson 2020 on media.

\notesec{Transcript against protein, and the two terms between them} Steady-state
transcript and steady-state protein are both well covered, by the served Kemmeren 2014 and
Messner 2023 and by the state rows, and they disagree. Two rows make the disagreement
something a model can carry a term for rather than absorb as error. McManus 2014 gives
translation efficiency beside mRNA abundance and finds ribosome occupancy more conserved
than transcript abundance, so translation buffers rather than tracks. Martin-Perez 2017
gives turnover for 3,160 proteins, spanning three orders of magnitude, which is why a
single global mRNA-to-protein scaling cannot be right. Sun 2013 supplies the transcript-side
counterpart, decomposing mRNA level into synthesis and decay across 46 deletion strains.

\notesec{Cross-laboratory replication of the deletion transcriptome} Kemmeren 2014 is
served and is the largest genotype-indexed transcriptome with designed perturbations, and
the served Hu 2007 now gives it one independent check on 263 transcription-factor
deletions. Hughes 2000, Lenstra 2011 and Sun 2013 add three more laboratories across
fourteen years of array and sequencing platforms, so agreement can be measured gene by
gene on shared strains rather than assumed. That check matters here specifically: the
served Nadal-Ribelles 2025 does not replicate its own genotype assignment across batches,
so a second opinion on what a deletion does to a transcriptome is not a formality.

\notesec{Metabolism against expression on one genotype axis} Blank 2005 measures flux
through central carbon metabolism on deletion mutants that already carry an expression
profile in the served Kemmeren 2014 compendium, which is the only place in this table
where the transcript-to-flux step can be fitted rather than assumed. Aulakh 2025 adds a
second metabolic panel, the ionome, on the whole collection. Gerke 2017, Ambroset 2014
and Eder 2020 put metabolite and flux quantitative trait loci on segregant crosses that
Albert 2018 and Jakobson 2025 give an expression and a protein map, so a metabolite locus
can be read against both.

\notesec{A condition axis measured on more than one layer} Brauer 2008 and Boer 2010 are
the transcript and metabolite halves of one chemostat series, same laboratory and same
limitations, and Hackett 2016 adds flux over the same conditions. Airoldi 2016 supplies the
nitrogen-limited counterpart, and its media are the media the served Jackson 2020 profiles
single cells in. That last pairing is what makes a deletion effect in the single-cell atlas
a contrast against a wild-type baseline in the same medium rather than an absolute.

%% GENERATED FILE -- do not hand-edit.
%% SOURCE: experiments/database/scripts/build_candidate_datasets_table.py
\begin{landscape}
\begingroup
\footnotesize
\setlength{\tabcolsep}{4pt}
\renewcommand{\arraystretch}{1.15}
\begin{longtable}{@{}r@{\hspace{4pt}} L{42mm} L{44mm} L{48mm} L{96mm}@{}}
\caption[]{Named joins between a candidate and either a supported dataset or
another candidate, for the sixty rows of Table~\ref{tab:final}. \emph{\#} is the
candidate's rank there. A partner in \textbf{bold} is already built, so that
row's join becomes runnable the moment the candidate is ingested; an unbolded
partner is itself a candidate, so the join costs two ingestions. \emph{Join key}
is the shared axis that makes the pair addressable; a pair with no such axis is a
theme and is not listed. None of these transfers has been measured.}
\label{tab:synergies}\\
\toprule
\textbf{\#} & \textbf{Candidate} & \textbf{Partner} & \textbf{Join key} & \textbf{What the join yields} \\
\midrule
\endfirsthead
\multicolumn{5}{@{}l}{\footnotesize\emph{Table~\ref{tab:synergies}, continued}}\\
\toprule
\textbf{\#} & \textbf{Candidate} & \textbf{Partner} & \textbf{Join key} & \textbf{What the join yields} \\
\midrule
\endhead
\bottomrule
\endfoot
1 & Albert 2018 (eQTL in 1,012 segregants) & Jakobson 2025 (genome-to-proteome map) & the same segregant panel & The transcript half of a paired transcript and protein map. \\
 &  & Gerke 2017 (urea-cycle mQTL) & segregant genotype class & A metabolite locus read through the transcripts of the pathway that produces it. \\
\addlinespace[5pt]
2 & Jakobson 2025 (genome-to-proteome map) & Albert 2018 (eQTL in 1,012 segregants) & the same sequenced segregant panel & Protein and transcript quantitative trait loci on one genotype set, which is the direct measurement of how much of protein variation transcript variation explains. \\
 &  & \textbf{Messner 2023 (proteome)} & protein abundance, natural variation against deletion & Whether the proteome response to a deleted gene resembles the response to a natural allele of it. \\
\addlinespace[5pt]
3 & Teyssonniere 2024 (species-wide proteome against transcriptome) & \textbf{Caudal 2024 (pan-transcriptome)} & the same sequenced isolates, 942 against 943 & Transcript and protein for the same strain, one to one rather than by overlap. The published result on this pair is that the two correlate weakly and share 3 percent of their associated variants, so it is the calibration an inference model is scored against. \\
 &  & \textbf{Muenzner 2024 (natural-isolate proteomes)} & the 1,011-isolate panel, protein layer & De-duplication before either is built: the two share authors and a panel, and whether they are independent acquisitions is unsettled. \\
\addlinespace[5pt]
4 & Grossbach 2022 (BY x RM transcriptome, proteome and phosphoproteome) & Albert 2018 (eQTL in 1,012 segregants) & BY x RM segregants, transcript layer & 112 strains with transcript and protein measured together against 1,012 with transcript alone, so the within-strain RNA-to-protein residual fitted on the small panel can be applied to the large one. \\
 &  & Jakobson 2025 (genome-to-proteome map) & segregant genotype class, protein layer & Two segregant proteomes on different panels and platforms, which is what separates a protein quantitative trait locus from a batch. \\
 &  & Leutert 2023 (phosphoproteome x 101 conditions) & phosphosite identity & Phosphorylation driven by genotype against phosphorylation driven by condition, on one site vocabulary. \\
\addlinespace[5pt]
6 & Skelly 2013 (expression variation across isolates) & \textbf{Muenzner 2024 (natural-isolate proteomes)} & natural isolates, paired against joined modalities & A within-experiment paired transcript and protein measurement to calibrate the cross-study join. \\
\addlinespace[5pt]
7 & McManus 2014 (ribosome profiling, allele-specific) & \textbf{Messner 2023 (proteome)} & gene identity, translation rate against protein abundance & A per-gene translation term to put against measured protein level, which is the coefficient a model otherwise has to learn blind. \\
 &  & Martin-Perez 2017 (protein half-lives) & gene identity, synthesis against degradation & Both halves of protein turnover, so steady-state abundance can be decomposed rather than only predicted. \\
\addlinespace[5pt]
8 & Foss 2007 (BY x RM segregant proteome) & Grossbach 2022 (BY x RM transcriptome, proteome and phosphoproteome) & the same cross, protein layer, fifteen years and two platforms apart & Cross-laboratory replication of segregant protein abundance, which neither measurement can establish alone. \\
\addlinespace[5pt]
9 & Dong 2021 (MAGIC + SAM biosensor) & \textbf{Lian 2019 (MAGIC CRISPR-AID)} & the same tri-functional library, fitness against biosensor readout & One library, two labels: growth selection and product concentration, which is how a tolerance screen is converted into a production screen. \\
 &  & \textbf{Cachera 2023 (CRI-SPA betaxanthin)} & metabolite biosensor or colorimetric product readout on a genome-scale library & Two product-labeled genome-scale screens on different chemistry. \\
\addlinespace[5pt]
10 & Momen-Roknabadi 2020 (inducible CRISPRi library) & McGlincy 2021 (genome-scale CRISPRi library) & independently designed guide libraries over the same genes & The paired half of the library-transfer test. \\
 &  & \textbf{SGD (essentiality)} & essential genes, knockdown against deletion & A graded phenotype where a deletion collection records only absence. \\
\addlinespace[5pt]
11 & McGlincy 2021 (genome-scale CRISPRi library) & Momen-Roknabadi 2020 (inducible CRISPRi library) & independently designed guide libraries over the same genes & Library-to-library transfer, which is the cheapest decisive test of whether a model learned gene function or guide-design idiosyncrasy. \\
 &  & \textbf{Smith 2016 (CRISPRi chem-genetic)} & CRISPR interference over the same gene set & A knockdown fitness prior for the chemical-genetic conditions already built. \\
\addlinespace[5pt]
12 & Bao 2018 (CHAnGE single-nucleotide library) & \textbf{Lian 2019 (MAGIC CRISPR-AID)} & guide-indexed genome-wide library from the same laboratory & Sub-gene alleles against whole-gene activation, interference and deletion on one platform lineage. \\
 &  & Li 2016 (tRNA fitness landscape) & designed single-nucleotide variants with a fitness label & Genome-wide shallow variant coverage against single-gene exhaustive coverage. \\
\addlinespace[5pt]
13 & Crook 2016 (tunable RNAi, isobutanol + 1-butanol) & \textbf{Lopez 2024 (isobutanol screen, private)} & isobutanol, knockdown against deletion & A graded knockdown axis over the phenotype the built biosensor screen measures as a titer proxy. \\
 &  & Kuroda 2019 (isobutanol-specific tolerance) & isobutanol tolerance on the same collection & Dose-graded knockdown against whole-gene deletion for one phenotype. \\
\addlinespace[5pt]
14 & Roy 2018 (multiplexed precision editing) & Bao 2018 (CHAnGE single-nucleotide library) & designed edits, single against multiplexed & Whether multiplexed edit effects are the sum of their single-edit effects. \\
\addlinespace[5pt]
15 & Newman 2006 (protein noise, GFP library) & \textbf{Messner 2023 (proteome)} & protein abundance across the genome, per cell against population & Which proteins have a population mean that no single cell is near, which decides where a mean-level model is the wrong object. \\
 &  & Gasch 2017 (single-cell stress heterogeneity) & per-cell variation in an isogenic population & Noise at the protein and transcript layers on the same question. \\
\addlinespace[5pt]
16 & Mukherjee 2021 (CRISPRi essential genes x acetic acid) & \textbf{Mormino 2022 (CRISPRi acetic-acid)} & acetic acid, CRISPR interference & The genome-scale extension of a twelve-gene screen already built. \\
 &  & \textbf{Mota 2024 (weak-acid screen)} & weak-acid stress on deletion strains & Essential-gene coverage for a phenotype the deletion collection can only sample from the non-essential side. \\
\addlinespace[5pt]
17 & Lian 2017 (CRISPR-AID, beta-carotene) & \textbf{Ozaydin 2013 (beta-carotene screen)} & beta-carotene titer & Three perturbation modalities against a whole-collection deletion screen on the same product. \\
\addlinespace[5pt]
18 & Hughes 2000 (compendium of expression profiles) & \textbf{Kemmeren 2014 (deletion compendium)} & deletion strains, gene by gene, fourteen years and two array platforms apart & Cross-laboratory replication of the deletion transcriptome, which the built set cannot test against itself and which Nadal-Ribelles 2025 failed cross-batch. \\
 &  & \textbf{Hu 2007 (TF-deletion error-model scores)} & deletion strains with a bulk expression readout & A third independent compendium, so agreement can be measured across three laboratories rather than asserted from two. \\
\addlinespace[5pt]
19 & Lenstra 2011 (chromatin regulator deletion expression) & \textbf{Kemmeren 2014 (deletion compendium)} & single deletions with a bulk expression readout, different gene class & Chromatin regulators beside sequence-specific factors on one expression readout. \\
 &  & \textbf{Sameith 2015 (dm)} & regulator deletions, single against double & Whether chromatin-regulator effects combine the way transcription-factor effects do. \\
\addlinespace[5pt]
20 & Sun 2013 (mRNA synthesis and decay rates across deletion strains) & \textbf{Kemmeren 2014 (deletion compendium)} & single deletion strains, gene by gene & A steady-state transcript level beside the synthesis and decay rates that produce it, so two strains with the same level but different kinetics stop looking identical. \\
 &  & Martin-Perez 2017 (protein half-lives) & turnover, transcript side against protein side, gene by gene & Both half-life terms in one place, which is what an RNA-to-protein model needs in order to have a reason for the two to disagree. \\
\addlinespace[5pt]
21 & Jariani 2020 (yeast scRNA-seq through lag phase) & \textbf{Jackson 2020 (TF-deletion pseudobulk)} & carbon-source shift with a per-cell readout & A transition sampled densely in time where the atlas samples eleven steady states. \\
\addlinespace[5pt]
22 & Guo 2018 (CRISPR-Cas9 tiling of essential genes) & \textbf{SGD (essentiality)} & essential genes & A graded allelic series where the built record is binary. \\
\addlinespace[5pt]
23 & Urbonaite 2021 (yeastDrop-Seq under drug treatment) & \textbf{Hoepfner 2014 (HIP/HOP atlas)} & chemical treatment of yeast & A per-cell readout for compound response where the built atlas has a pooled fitness score. \\
\addlinespace[5pt]
24 & Nadal-Ribelles 2019 (sensitive yeast scRNA-seq) & \textbf{Nadal-Ribelles 2025 (Perturb-seq)} & same laboratory, protocol against application & The sensitivity ceiling of the platform the built genome-scale single-cell dataset was produced on. \\
 &  & Gasch 2017 (single-cell stress heterogeneity) & isogenic per-cell variation in BY4741 & Two independent measurements of how much an unperturbed yeast population varies, which is the floor a perturbation must clear. \\
\addlinespace[5pt]
25 & Gasch 2017 (single-cell stress heterogeneity) & Gasch 2000 (environmental stress response) & the same stress conditions, bulk against single cell & How much of the canonical stress response is a population average of cells that are individually in different states. \\
\addlinespace[5pt]
26 & Jaffe 2019 (multiplexed CRISPR interference epistasis) & \textbf{Costanzo 2016 (dmi)} & gene pairs, knockdown against deletion & Whether digenic interaction measured between two nulls is recovered between two knockdowns, which decides if the built interaction data can supervise a knockdown campaign. \\
 &  & \textbf{Kuzmin 2018 (tmi)} & higher-order gene combinations & The knockdown analog of trigenic interaction. \\
\addlinespace[5pt]
27 & Jackson 2023 (wild-type scRNA-seq time course for RNA kinetics) & \textbf{Jackson 2020 (TF-deletion pseudobulk)} & same platform and laboratory, wild type against deletions & The unperturbed per-cell variance the perturbed atlas is read against. \\
\addlinespace[5pt]
28 & Brettner 2024 (ultra-high-throughput yeast scRNA-seq) & \textbf{Jackson 2020 (TF-deletion pseudobulk)} & multiplexed genotypes per single-cell run & The throughput ceiling for a barcoded-genotype design, measured rather than assumed. \\
\addlinespace[5pt]
29 & Mulleder 2012 (prototrophic deletion collection) & \textbf{Mulleder 2016 (amino-acid metabolome)} & the same prototrophic deletion strains & The strain resource the built amino-acid metabolome was measured on, which is what makes defined-medium phenotyping interpretable. \\
 &  & \textbf{Jackson 2020 (TF-deletion pseudobulk)} & prototrophy as a precondition for minimal media & A ready-made prototrophic library for a campaign that needs nitrogen or carbon limitation. \\
\addlinespace[5pt]
30 & Su 2023 (single-cell transcriptomes under four stresses) & Gasch 2000 (environmental stress response) & osmotic and starvation stress & Full-length per-cell transcripts against the bulk stress-response reference. \\
\addlinespace[5pt]
31 & Wang 2022 (single-cell transcriptomes across replicative aging) & Jackson 2023 (wild-type scRNA-seq time course for RNA kinetics) & wild-type per-cell states over time & Replicative age as a covariate on a platform that does not measure it. \\
\addlinespace[5pt]
32 & Aulakh 2025 (genome-scale ionome) & \textbf{Mulleder 2016 (amino-acid metabolome)} & the whole deletion collection, metabolite readout & Two orthogonal metabolic panels on one genotype axis. \\
 &  & \textbf{Kemmeren 2014 (deletion compendium)} & the same deletion collection & Element levels against the transcriptional response of the same knockout. \\
\addlinespace[5pt]
33 & Leutert 2023 (phosphoproteome x 101 conditions) & Gasch 2000 (environmental stress response) & overlapping stress conditions & Post-translational regulation against transcriptional regulation for one condition set, on the layer that acts first. \\
 &  & \textbf{Messner 2023 (proteome)} & protein identity & Modification state against abundance, which abundance alone cannot separate. \\
\addlinespace[5pt]
34 & Brauer 2008 (growth-rate-controlled chemostat transcriptome) & Boer 2010 (metabolome across nutrient limitations) & the same chemostat limitations & The transcript half of a paired transcript and metabolite series. \\
 &  & Airoldi 2016 (nitrogen-limited steady-state and dynamic transcriptome) & growth rate under limitation, carbon against nitrogen & Whether the growth-rate expression program is nutrient-general. \\
\addlinespace[5pt]
35 & Hackett 2016 (SIMMER multi-omic flux) & Boer 2010 (metabolome across nutrient limitations) & the same nutrient-limited chemostats & Metabolite concentration and flux under one condition set, which is what a kinetic model needs and neither supplies alone. \\
 &  & Brauer 2008 (growth-rate-controlled chemostat transcriptome) & nutrient limitation at controlled growth rate & Expression, metabolite and flux over one condition axis. \\
\addlinespace[5pt]
36 & Zhu 2014 (kinase / phosphatase lipidomics) & \textbf{da Silveira 2014 (lipidomics)} & lipid species on deletion strains & A second lipidome on a signaling-focused genotype set. \\
 &  & \textbf{Kemmeren 2014 (deletion compendium)} & kinase and phosphatase deletions & Lipid composition against the transcriptional response of the same signaling mutant. \\
 &  & \textbf{Xue 2025 (free fatty acids, private)} & fatty-acid metabolism & Lipid class distribution against free fatty-acid titer, the two sides of the malonyl-CoA sink. \\
\addlinespace[5pt]
37 & Gerke 2017 (urea-cycle mQTL) & Albert 2018 (eQTL in 1,012 segregants) & segregant genotype class & A metabolite locus and the expression loci in the same interval. \\
 &  & \textbf{Mulleder 2016 (amino-acid metabolome)} & nitrogen and amino-acid metabolites & Natural-allele against deletion effects on one metabolic module. \\
\addlinespace[5pt]
38 & Ambroset 2014 (metabolite QTL) & Peeters 2021 (fermentation-trait QTL atlas) & segregant panels with fermentation traits & 74 metabolites against 18 mapped traits, so a trait locus can be read as a metabolite change. \\
 &  & \textbf{Zelezniak 2018 (metabolome)} & metabolite panel & Recombinant against deletion genotypes on a metabolite readout. \\
\addlinespace[5pt]
39 & Blank 2005 (13C metabolic flux) & \textbf{Kemmeren 2014 (deletion compendium)} & deletion strains present in both & Flux against expression for the same knockout, which is the only place the transcript-to-flux step can be fitted rather than assumed. \\
 &  & \textbf{Mulleder 2016 (amino-acid metabolome)} & deletion strains, flux against pool size & Whether a changed pool reflects changed flux or changed demand. \\
\addlinespace[5pt]
40 & Airoldi 2016 (nitrogen-limited steady-state and dynamic transcriptome) & \textbf{Jackson 2020 (TF-deletion pseudobulk)} & the same nitrogen-limited media & A wild-type bulk baseline for the medium the single-cell atlas perturbs in. \\
 &  & Yu 2021 (proteome and metabolome under nitrogen limitation) & nitrogen limitation & Transcript, protein and metabolite under one limitation. \\
\addlinespace[5pt]
41 & Martin-Perez 2017 (protein half-lives) & \textbf{Messner 2023 (proteome)} & gene identity, half-life against abundance & Which proteins are abundant because they are made fast and which because they are destroyed slowly, a distinction abundance alone cannot make. \\
 &  & \textbf{Kemmeren 2014 (deletion compendium)} & gene identity, protein half-life against transcript response & Whether a transcript change reaches the protein layer at all, which for a long-lived protein it largely does not. \\
\addlinespace[5pt]
42 & Boer 2010 (metabolome across nutrient limitations) & Brauer 2008 (growth-rate-controlled chemostat transcriptome) & the same chemostat limitations and growth rates & The metabolite and transcript halves of one experimental series. \\
 &  & \textbf{Zelezniak 2018 (metabolome)} & metabolite panel & Condition-driven against genotype-driven metabolite variation. \\
\addlinespace[5pt]
43 & Tengolics 2024 (domestication metabolome) & \textbf{Caudal 2024 (pan-transcriptome)} & sequenced natural isolates & Metabolite against transcript variation across the same population. \\
\addlinespace[5pt]
44 & Eder 2020 (flux QTL) & Blank 2005 (13C metabolic flux) & flux measurement, natural variation against deletion & Whether flux control points found by deletion are the loci natural variation actually moves. \\
 &  & Albert 2018 (eQTL in 1,012 segregants) & segregant genotype class & Flux loci against expression loci on comparable panels. \\
\addlinespace[5pt]
45 & Yu 2021 (proteome and metabolome under nitrogen limitation) & Airoldi 2016 (nitrogen-limited steady-state and dynamic transcriptome) & nitrogen limitation & The transcript layer for the protein and metabolite layers measured here. \\
\addlinespace[5pt]
46 & de Boer 2020 (100M random promoters) & \textbf{Vaishnav 2022 (promoter sort-seq)} & the same random promoter library, expression against fitness & Two labels on one sequence set, which turns a sequence-to-expression model into a sequence-to-fitness model without new data. \\
 &  & Renganaath 2020 (natural promoter-variant MPRA) & promoter sequence with a measured expression readout & Whether a model trained on random sequence predicts the effect of a real allele, the cheapest decisive transfer test in the table. \\
\addlinespace[5pt]
47 & Lee 2014 (HIP-HOP fitness signatures) & \textbf{Hoepfner 2014 (HIP/HOP atlas)} & the same heterozygous and homozygous collections, compound response & Two chemical-genetic matrices on one genotype axis, which is enough compound overlap to measure cross-screen reproducibility directly. \\
 &  & \textbf{Hillenmeyer 2008 (FitDb HIP, het)} & the same collections and readout & A third independent screen of the same design. \\
\addlinespace[5pt]
48 & Nguyen Ba 2022 (barcoded bulk QTL, 100k segregants) & \textbf{Bloom 2019 (16-cross segregant panel)} & segregant genotype class, panel size & One cross at 100,000 progeny against sixteen founders at 14,000, which separates panel size from allelic diversity. \\
\addlinespace[5pt]
49 & Galardini 2019 (four backgrounds x 38 conditions) & \textbf{Costanzo 2021 (condition-SGA)} & deletion strains across conditions & The same conditional fitness question asked in four backgrounds rather than one. \\
 &  & \textbf{Hale 2024 (CRISPRi x segregants)} & perturbation held fixed, background varied & Deletion against knockdown for the background-dependence question. \\
\addlinespace[5pt]
50 & Dutta 2026 (barcoded natural-isolate chemical response) & \textbf{Caudal 2024 (pan-transcriptome)} & the sequenced 1,011-isolate panel & Chemical response and transcriptome on one isolate set. \\
 &  & \textbf{Hoepfner 2014 (HIP/HOP atlas)} & compound response, isolates against deletions & Whether compound sensitivity found by gene deletion predicts sensitivity across natural genetic backgrounds. \\
\addlinespace[5pt]
51 & Li 2016 (tRNA fitness landscape) & Domingo 2018 (tRNA double-mutant landscape) & the same tRNA gene, single against double mutants & Whether a single-mutant landscape predicts the double-mutant one, which is the within-gene form of the epistasis question. \\
\addlinespace[5pt]
58 & Liu 2021 (tryptophan / isobutanol tolerance) & Kuroda 2019 (isobutanol-specific tolerance) & isobutanol tolerance on the same deletion collection & Same genotypes, same phenotype, different readout, which is a method-effect control no single screen provides. \\
\addlinespace[5pt]
59 & Kuroda 2019 (isobutanol-specific tolerance) & \textbf{Lopez 2024 (isobutanol screen, private)} & isobutanol on the deletion collection & Tolerance against production for one product, on one genotype axis. \\
\addlinespace[5pt]
60 & Turco 2023 (Yeast Phenome) & \textbf{Hoepfner 2014 (HIP/HOP atlas)} & aggregated growth screens against a primary screen & De-duplication, which is the precondition for knowing what the aggregate adds. \\
\addlinespace[5pt]
\end{longtable}
\endgroup
\end{landscape}

